\begin{proposition}[The period-conjecture implication for the local norm]\label{ari:C.7.4}
Let $E/\overline\Q$ be an elliptic curve and let
$P\in E(\overline\Q)$ be nontorsion. Use the local norm $N(P)$,
the algebraic extension $G_P$, and its algebraic fibre coordinate
from \cref{ari:C.7.1,ari:C.7.2}.
The hypothesis $H$ is the Relative Semi-abelian Conjecture
\cite[Conjecture~4.7, p.~18]{bertolin2022}, restricted to
$G=G_P\times G_{\bar P}$ and the algebraic lifts used below.
Explicitly, for those lifts and their associated $B_Q,Z(1)$, it gives
\[
 \operatorname{trdeg}_{\overline\Q(\Omega_G)}
 \overline\Q\bigl(\Omega_G,R,\bar R,
                  \widetilde{\log}_G(R,\bar R)\bigr)
 \geq 2\dim B_Q+\dim Z(1),
\]
where $\Omega_G$ is the full de Rham--Betti period matrix,
including the second-kind periods. Under $H$, $N(P)$ is
transcendental. This applies to every algebraic parameter
$t\ne0,1$ with nontorsion $P_t$, including the conductor-$361$
CM point identified below. Hypothesis $H$ remains open for
these points. Grothendieck's period conjecture for the algebraic
$1$-motives in this proof implies $H$
\cite[Proposition~4.12, pp.~19--20]{bertolin2022}.
\end{proposition}
\Uses{\cref{ari:C.7.2,ari:C.7.3,ari:eq:cp-v,ari:eq:cp-positive-centre};
\cite[Definitions~2.3--2.4, p.~8, Construction~3.4, p.~12,
Lemma~3.5, p.~13, Conjecture~4.7, p.~18,
Proposition~4.12 and Remark~4.13, pp.~19--20]{bertolin2022}}
\begin{proof}
\begin{enumerate}[label=\arabic*.,leftmargin=*]
\item Suppose that $N(P)$ is algebraic. The compact-character
identity of \cref{ari:C.7.2} supplies an algebraic point
$R=(P,r)\in G_P(\C)^c$: we take the positive algebraic number
$r=|\wp'(u)|\,|\Delta|^{-1/6}N(P)^2$ in the specified
algebraic fibre coordinate. The conjugate point gives the
corresponding algebraic lift $\bar R$ over $\bar P$.
Set $M_R=[\Z\to G_P]$, $1\mapsto R$, and
$\mathcal M=M_R\oplus\overline{M_R}$.

\item Put $K_0=\overline\Q(\Omega_G)$. Write
$P=[u]$ with periods $\omega_1,\omega_2$, quasi-periods
$\eta_1,\eta_2$, and the logarithmic coordinate $v$ of
\eqref{ari:eq:cp-v}. The full matrix contains
\[
 \kappa_j=\zeta(u)\omega_j-\eta_j u,\qquad j=1,2.
\]
The coefficient matrix for $(\zeta(u),u)$ is
$\left(\begin{smallmatrix}\omega_1&-\eta_1\\
\omega_2&-\eta_2\end{smallmatrix}\right)$; its determinant
is the nonzero Legendre determinant. Hence $u,\zeta(u)$
belong to $K_0$. The same calculation in the conjugate factor
places their conjugates there. This is the explicit elliptic
instance of \cite[Remark~4.13, p.~20]{bertolin2022}.

\item Write $u=a\omega_1+b\omega_2$ with $a,b\in\R$.
Solving this system together with its conjugate gives
\begin{equation}\label{ari:eq:cp-real-linear}
 \eta(u)=
 \frac{(u\bar\omega_2-\bar u\omega_2)\eta_1
       +(\bar u\omega_1-u\bar\omega_1)\eta_2}
      {\omega_1\bar\omega_2-\bar\omega_1\omega_2}.
\end{equation}
The denominator is nonzero by real-linear independence of the
periods. Thus $u\eta(u)$ and its conjugate belong to $K_0$.
The compact-character equation is
$\Re(v-u\zeta(u)+u\eta(u))=0$, so
\begin{equation}\label{ari:eq:cp-relation}
 v+\bar v=u\zeta(u)+\bar u\overline{\zeta(u)}
                -u\eta(u)-\bar u\overline{\eta(u)}\in K_0.
\end{equation}
Each conjugate is a coordinate of the conjugate algebraic
$1$-motive, whose full periods were included in $K_0$.

\item The period torsor of $\mathcal M$ maps to the period
torsor of $[0\to G]$. Projection of the abelian point-extension
variety $B$ to the extension-coordinate variety is an isogeny:
the point coordinates are the same algebraic data as the
extension coordinates under the specified polarizations.
Equivalently, in the product $E\times\bar E$, the projected
point $(P,\bar P)$ is the sum of the two extension parameters
$(P,0)$ and $(0,\bar P)$. The formula for $B_Q$ in
\cite[equation~(4.4), p.~19]{bertolin2022} gives $\dim B_Q=0$.
There is therefore no weight-$-1$ kernel. The full kernel on
the connected motivic groups is the weight-$-2$ vector group
$\operatorname{Lie}Z$, as in
\cite[Construction~3.4 and Lemma~3.5, pp.~12--13]{bertolin2022}.
After fixing the base periods its fibre is an affine torsor
over $K_0$ of dimension $\dim Z$.

\item In the positive Tate basis of each algebraic factor,
the central action changes $(v,\bar v)$ by
$2\pi i(c_1,c_2)$. Replacing a conjugated negative Betti loop
by the positive loop changes its period column, while leaving
the analytic coordinate $\bar v$ fixed. The polarization
transport through $\phi:\bar E\to E$ of degree $d$ is
$\phi^tJ_E\phi=dJ_{\bar E}$. Thus the bracket line computed
in \cref{ari:C.7.3} has coordinates
$(1,aa^\dagger/d)$ with $aa^\dagger/d>0$; in the independent
case the bracket fills both central coordinates. In both cases
$(c_1,c_2)\mapsto c_1+c_2$ is nonzero on the actual central
group. Any extra quotient Tate direction preserves this
nonzero restriction. Consequently \eqref{ari:eq:cp-relation}
cuts a proper affine hyperplane in that torsor and forces the
relative transcendence degree to be at most $\dim Z-1$.
The same doubled construction applies to real data; its
rank-one bracket still has positive coefficient.

\item All noncentral generalized logarithm coordinates are in
$K_0$ by steps 2--3, and $R,\bar R$ are algebraic. Hypothesis
$H$ gives the lower bound $\dim Z$, contradicting step 5.
This proves the stated implication. The source's
Proposition~4.12 supplies it from the period conjecture for
these algebraic $1$-motives. No assertion of $H$ is used
elsewhere to remove its stated open status.
\end{enumerate}
\end{proof}
